\subsection{Case study data}
The experiments use two real tunnel cases. The first is the BSLL tunnel case
starting from DyK1017+205. It provides the main graph-sequence construction
example, the one-step response prediction run, and the three-step
advance-prediction sensitivity run. The second is the SJLS tunnel case starting
from Dyk1252+411. It uses an external TSP-derived velocity field and real
mileage-aligned TBM monitoring labels, and is used to test whether the same
spatial representation remains meaningful in a second tunnel setting.

The two cases are not treated as sequential data releases or replacement
datasets. They are two case-study settings with different evidential roles.
BSLL is used mainly to illustrate multi-step prediction and spatial
interpretation under a compact test interval. SJLS is used mainly for
prediction consistency and geometry-ablation evidence under an external TSP
case.

\begin{figure}[htbp]
\centering
\includegraphics[width=0.9\linewidth]{fig2_graph_construction.pdf}
\caption{Schematic of geometry-constrained graph construction at a
representative excavation step (illustrative, not drawn from the actual TSP
velocity field). The framework links TSP-derived rock voxels, parameterized TBM
surface nodes, and geometry-screened rock--machine candidate edges before
learning response-consistent relevance.}
\label{fig:tsp_profile}
\end{figure}

\subsection{Monitoring and geological variables}
For each excavation step, the monitoring vector contains advance rate, cutterhead
RPM, torque, thrust, penetration, and shield pressure. The prediction target is
the future multivariate response vector consisting of advance rate, torque,
thrust, penetration, and shield pressure. These variables describe the observed
machine response to the surrounding rock and are used both as historical inputs
and as response-supervision targets.

The geological input is a TSP-derived voxel field. For BSLL, the project TSP
voxel field contains P-wave velocity, S-wave velocity, density, dynamic Young's
modulus, velocity ratio, and Poisson's ratio around the tunnel alignment. For
SJLS, dense external P-wave and S-wave velocity profiles at Dyk1252+411 are
preprocessed into the same project TSP schema and paired with mileage-aligned
monitoring labels. This preprocessing allows both tunnel cases to use the same
graph-construction and graph-sequence learning pipeline.

\subsection{Mileage-ordered partitioning}
All experiments use mileage-ordered partitioning to preserve forward prediction.
For a historical window of $K=5$, each sample uses the preceding monitoring and
graph snapshots as input and predicts a future response at horizon $h$. The
formal case runs are summarised in Table~\ref{tab:data_summary}. The BSLL h=1
run evaluates one-step prediction, the BSLL h=3 run evaluates multi-step
advance prediction, and the SJLS run evaluates the external TSP case.

\begin{table}[htbp]
\centering
\caption{Formal case runs used in the experiments. Effective samples are
sequence windows after applying the history window and prediction horizon.}
\label{tab:data_summary}
\begin{tabular}{lcccc}
\toprule
Case & Tunnel & Start chainage & Horizon & Train/Test samples \\
\midrule
BSLL h=1 & BSLL & DyK1017+205 & 1 & 30 / 8 \\
BSLL h=3 & BSLL & DyK1017+205 & 3 & 29 / 7 \\
SJLS h=1 & SJLS & Dyk1252+411 & 1 & 76 / 17 \\
\bottomrule
\end{tabular}
\end{table}
